\subsection{我们该如何建立良性循环？}

短答案是：让下一次好的动作比下一次坏的动作更容易发生，然后不断重复这种优势，直到新模式开始自己带动自己。
良性循环不是靠一次大爆发的意志力建立起来的。
它是通过安排环境、时机和第一步，让每一次成功的转变都略微降低下一次障碍、并稍微增强目标状态而形成的。
用本书的语言说，目标是在进入时刻降低框架惯性，把意志能量用在杠杆最大的地方，并让决策窗口短到旧框架来不及完全夺回控制。

\begin{enumerate}[nosep]
  \item 用朴素语言命名一个具体目标状态。不要先从身份宣言开始，要先从你想让它变容易的行为开始。
  \item 找出那个最小的、仍能完成转变的决策窗口。如果窗口太长，就用预先规定第一步的方式把它缩短。
  \item 在下一次尝试前移除一个障碍变量。移动物件、预备工作台、降低模糊性，或者消除一个可预测的打断源。
  \item 规定一个最小可行进入动作，小到即使意志能量只有中等水平也能启动。
  \item 连续数次重复同样的进入条件，让目标状态开始显得熟悉，从而变得不那么昂贵。
  \item 每次成功之后，都做一个能够降低下一次启动成本的局部改进，而不是只追加更宏大的愿望。
\end{enumerate}

\begin{remark}
一个想建立学习良性循环的学生，不该先承诺“以后要更自律”。
更好的做法，是把书提前翻在桌上，先写下第一题的题号，并规定这次会话先做五分钟。
第一次成功会降低第二次障碍，而常规会变容易，是因为系统开始预期那个目标状态会发生。
\end{remark}